\subsection{Measured costs of the certified loop removals}

The two identity collapses above were applied at their particular
source occurrences. Shared occurrences elsewhere were retained; a
global replacement would lack the required zero premise. The independent
checker verifies all feasible incoming and outgoing edges through
normalization, not merely the local tree rewrite.

\begin{center}
\small\begin{tabular}{lrrrr}
\toprule
Removed zero loop & Source leaves & Retained cases & Emitted states & Best after one rewrite\\
\midrule
$X$ loop &6,457&8,384&253&253\\
$T$ loop &6,457&8,384&247&247\\
Both &6,454&8,383&254&254\\
\bottomrule
\end{tabular}
\end{center}

Each collapse preserves the original root and skips at most one
reachable identity source step between retained leaves. The two
additional deleted body leaves of each loop have empty invariant.
The three independent source certificates therefore prove the stated
source simplifications on the inherited product invariant.

The cost experiment checked all 718 closed one-step rewrite proposals
across these three starting programs, as well as the three unrewritten
programs. Unlike the earlier plateau cost gate, this final screen
computed exact one-step state bisimulation for \emph{every} emitted
graph, including those with more than 249 states. None improved its
respective starting count. The best source specialization costs 247
states under the current emitter, so no new complete binary candidate
results from this experiment.

For the closest case, the $T$-loop collapse, a further 333 one-step
control-aware global reassociations were checked. Exactly 257 passed
the invariant-conditioned source-edge repair and 76 failed it. Exact
bisimulation of the emitted tables still gave a minimum of 247.
This is again a one-step source and cost screen, not a universal binary
correctness or initialization certificate for those tables.

This is a concrete counterexample to the heuristic that removing
source instructions necessarily reduces the state table. Specializing
one occurrence destroys some shared decoder shapes and requires new
ancestors. The source proof succeeds; the measured compression goal
does not. A future compiler may recover the lost sharing through a
different representation, but these data do not establish such a
compiler or a saving below 246.
